% Irodov Part III §3.6, original problems 289--304.
% Source uploaded PDF printed pp.135--138 (PDF pp.144--147); answers pp.307--308.
\electroinductionproblem{289}{135}
图3.79中，均匀磁场 $\bm B$ 垂直矩形回路平面，长为 $\ell$ 的活动导体杆以速度 $v$ 平移，杆电阻为 $R$，左右两边导线 AB、CD 的电阻分别为 $R_1,R_2$。忽略自感，求活动杆中的电流。
\exsolution
运动杆上的动生电动势为 $\mathcal E=Bv\ell$。从活动杆的两端看，AB、CD 两条静止支路并联，其等效电阻为 $R_\parallel=R_1R_2/(R_1+R_2)$。因此杆的总电流
\[
\boxed{I_{\rm rod}=\frac{Bv\ell}{R+\dfrac{R_1R_2}{R_1+R_2}}}.
\]
两静止支路的分流遵守 $I_1/I_2=R_2/R_1$。
\excheck $R_1,R_2\to\infty$ 时回路开路，电流趋零；$R=0$ 时电流由并联负载限制。方向由 $\bm v\times\bm B$ 在杆上的投影确定。

\electroinductionproblem{290}{135}
半径 $a=25\,\mathrm{cm}$ 的金属圆盘绕轴匀速转动，角速度 $\omega=130\,\mathrm{rad/s}$。求中心与边缘的电势差：(a) 无外加磁场；(b) 有垂直圆盘的均匀磁场 $B=5.0\,\mathrm{mT}$。
\exsolution
电子质量 $m_e$、电荷模 $e$。在无磁场时，稳态径向电场需平衡电子所需向心加速度，故电势差模
\[
\boxed{|\Delta\varphi_a|=\frac{m_e\omega^2a^2}{2e}
\simeq3.0\,\mathrm{nV}}.
\]
(b) 磁场中圆盘半径 $r$ 处速度 $\omega r$，洛伦兹力的径向分量由静电场平衡，主要电场模 $E_r\simeq\omega rB$，从而
\[
\boxed{|\Delta\varphi_b|\simeq
\int_0^a\omega rB\,\mathrm dr
=\frac{\omega Ba^2}{2}\simeq20\,\mathrm{mV}}.
\]
\excheck (b) 中电子惯性对应的 nV 量级远小于磁动生电压的 mV 量级；电压符号需另按 $\bm\omega$ 与 $\bm B$ 的相对方向确定。

\electroinductionproblem{291}{135}
直径为 $d=20\,\mathrm{cm}$ 的半圆形细导线 AC 绕经过端点 A、且垂直直径 AC 的轴以 $\omega=100\,\mathrm{rad/s}$ 转动，均匀磁场 $B=5.0\,\mathrm{mT}$，$\bm\omega$ 与 $\bm B$ 同向。求沿导线 A 至 C 的 $\int_A^C\bm E_{\rm eff}\cdot\mathrm d\bm r$，并推广至其它形状的连接导线。
\exsolution
匀速转动导线上一点 $\bm r$ 的速度为 $\bm v=\bm\omega\times\bm r$；导体有效驱动力场为 $\bm v\times\bm B$。当 $\bm\omega\parallel\bm B$ 时，沿离转轴垂距 $r$ 的径向分量为 $\omega Br$。因此线积分只依赖两个端点到转轴的距离：
\[
\int_A^C(\bm v\times\bm B)\cdot\mathrm d\bm\ell
=\frac{\omega B}{2}(r_C^2-r_A^2).
\]
按原图从 A 至 C 的正向与场方向相反，$r_A=0,\ r_C=d$，所以
\[
\boxed{\int_A^C\bm E_{\rm eff}\cdot\mathrm d\bm r
=-\frac12\omega Bd^2=-10\,\mathrm{mV}}.
\]
\excheck 路径可以弯曲，结果只取决于端点到转轴的距离及积分方向；$\bm E_{\rm eff}$ 指运动电动势场，不能将其与单纯静电保守场混淆。

\electroinductionproblem{292}{135--136}
图3.80中半径 $a$ 的半圆形回路位于均匀磁场半空间的边界，磁场垂直图面。回路自 $t=0$ 起绕边界处的 O 轴以恒角加速度 $\beta$ 转动。求感应电动势 $\mathcal E(t)$ 并说明图形。
\exsolution
角位移为 $\theta(t)=\beta t^2/2$，角速度 $\dot\theta=\beta t$。每转过半周，半圆回路与有磁场的半空间重叠的面积随转角在 $0$ 和 $\pi a^2/2$ 之间线性变化，故除边界瞬间外
$|\mathrm dS_{\rm overlap}/\mathrm d\theta|=a^2/2$。由法拉第定律
\[
|\mathcal E|=B\left|\frac{\mathrm dS_{\rm overlap}}{\mathrm dt}\right|
=\boxed{\frac12Ba^2\beta t}.
\]
按图给定正向，令 $n=1,2,\ldots$ 表示当前第 $n$ 个半周，则
\[
\boxed{\mathcal E(t)=\frac12(-1)^n Ba^2\beta t}.
\]
符号切换时刻由 $\beta t_n^2/2=n\pi$ 给出：
$\boxed{t_n=\sqrt{2\pi n/\beta}}$。
\excheck 图线是模随时间线性增大、符号在每个半周转换处跳变的交替锯齿形；$Ba^2\beta t$ 的单位为 V。

\electroinductionproblem{293}{136}
图3.81的无限长直线电流 $I$ 与回字形导轨共面，导轨中长度 $\ell$、电阻 $R$ 的活动杆以速度 $v$ 远离直导线，杆距直导线 $r$。忽略导轨电阻和自感，求感应电流。
\exsolution
直导线在距轴 $r$ 的磁感应强度
$B(r)=\mu_0I/(2\pi r)$，垂直回路平面。活动杆位于固定距离 $r$，整根杆运动电动势
$\mathcal E=B(r)\ell v$，因此
\[
\boxed{I_{\rm ind}(r)=\frac{\mu_0I\ell v}{2\pi Rr}}.
\]
由楞次定律，随杆右移，原直导线贯穿回路的磁通量减少，感生电流必须产生同向磁场；方向由原图中直导线电流向上时的右手定则确定。
\excheck $I_{\rm ind}\propto1/r$，不是 $1/r^2$，因为动生电动势取活动杆所在位置的局部磁场。

\electroinductionproblem{294}{136}
图3.82边长为 $a$ 的正方形线框与长直电流 $I$ 共面，线框以速度 $v$ 远离直导线，其靠近直导线的边与导线距离为 $x$。求线框感应电动势。
\exsolution
穿过线框的磁通量（以固定法向计）
\[
\Phi(x)=\int_x^{x+a}\frac{\mu_0I}{2\pi r}a\,\mathrm dr
=\frac{\mu_0Ia}{2\pi}\ln\frac{x+a}{x}.
\]
因此
\[
\mathcal E=-\frac{\mathrm d\Phi}{\mathrm dt}
=-\frac{\mu_0Iav}{2\pi}
\left(\frac1{x+a}-\frac1x\right).
\]
模为
\[
\boxed{|\mathcal E|=\frac{\mu_0Ia^2v}{2\pi x(x+a)}}.
\]
\excheck 书末答案的 PDF 文本抽取在 $a^2$ 位置可能漏掉平方，依据磁通积分和单位 $[\mathcal E]=\mathrm V$ 必须为 $a^2$。

\electroinductionproblem{295}{136}
图3.83中质量为 $m$、长度为 $a$ 的金属杆可绕水平轴 O 在半径为 $a$ 的圆形导轨上滑动，垂直纸面的恒定磁场为 $B$，回路电阻 $R$。求为使杆以恒角速度 $\omega$ 转动，所需电源电动势随时间的变化；忽略摩擦、自感与轨道电阻。
\exsolution
转动杆的动生电动势为
$\mathcal E_{\rm mot}=\int_0^aB\omega r\,\mathrm dr=B\omega a^2/2$。杆电流的磁力矩为
$\tau_m=\int_0^aIBr\,\mathrm dr=IBa^2/2$。
杆的重力矩按图示转角 $\omega t$ 为
$\tau_g=(mga/2)\sin\omega t$。
恒角速度意味着净力矩为零，故所需电流
$I=mg\sin\omega t/(Ba)$。
回路电源应补偿电阻电压降和动生电动势：
\[
\boxed{\mathcal E(t)=\frac{B\omega a^2}{2}
+\frac{mgR}{Ba}\sin\omega t
=\frac{B^2\omega a^2+2mgR\sin\omega t}{2Ba}}.
\]
\excheck 与原书答案第295题一致；电动势可能在某些相位反号，表示电源需要反向吸收能量。

\electroinductionproblem{296}{136--137}
图3.84中质量为 $m$ 的铜质活动杆沿与水平夹角为 $\alpha$ 的两根光滑导轨下滑，导轨间距为 $\ell$，顶部用电阻 $R$ 连接，均匀磁场 $B$ 垂直导轨平面。忽略其它电阻及自感，求最终稳定速度。
\exsolution
沿导轨运动速度 $v$ 产生电动势
$\mathcal E=B\ell v$，回路电流 $I=B\ell v/R$。磁阻力
$F_m=BI\ell=B^2\ell^2v/R$，沿斜面向上。稳定时
$mg\sin\alpha=F_m$，故
\[
\boxed{v_\infty=\frac{mgR\sin\alpha}{B^2\ell^2}}.
\]
\excheck 该结果与斜面加速度的终态 $0$ 相符；$R\to0$ 时强磁阻尼使终速趋零。

\electroinductionproblem{297}{137}
将第296题的顶部电阻 $R$ 换成电容 $C$，其余条件相同。求活动杆下滑的加速度。
\exsolution
活动杆的动生电动势为
$V_C=B\ell v$。电容电荷
$Q=CV_C=CB\ell v$，所以回路瞬时电流
$I=\dot Q=CB\ell a$。该电流在活动杆上产生向上磁阻力
$F_m=B\ell I=CB^2\ell^2a$。
沿导轨方向应用牛顿第二定律：
$ma=mg\sin\alpha-CB^2\ell^2a$。
解得
\[
\boxed{a=\frac{g\sin\alpha}{1+CB^2\ell^2/m}}.
\]
\excheck 电容产生的力与加速度而不是速度成正比；由量纲 $CB^2\ell^2$ 是质量。

\electroinductionproblem{298}{137}
图3.85半径为 $a$ 的半圆形导线绕轴 OO' 以角速度 $\omega$ 在与转轴垂直的均匀磁场 $B$ 中旋转，整个闭合回路电阻为 $R$。忽略感应电流的磁场，求一周期内平均焦耳热功率。
\exsolution
半圆所围面积为 $S=\pi a^2/2$，绕直径轴旋转使磁通为
$\Phi(t)=BS\cos\omega t$（取合适初相）。法拉第电动势
$\mathcal E=BS\omega\sin\omega t$。
于是
\[
\overline P=\frac{\overline{\mathcal E^2}}R
=\frac{B^2S^2\omega^2}{2R}
=\boxed{\frac{\pi^2a^4B^2\omega^2}{8R}}.
\]
\excheck $\overline{\sin^2\omega t}=1/2$；结果依赖半圆面积平方，而非整圆面积。

\electroinductionproblem{299}{137}
小线圈截面积 $S=3.0\,\mathrm{mm^2}$、匝数 $N=60$，初始轴线平行于磁极间均匀磁场，将线圈绕其直径翻转 $180^\circ$，与其相连的冲击电流计累计电荷为 $q=4.5\,\mathrm{\mu C}$。总回路电阻 $R=40\,\Omega$，求 $B$。
\exsolution
翻转前磁链为 $+NBS$，翻转后为 $-NBS$，磁链变化量大小 $2NBS$。对感应电路积分：
$q=\int I\,\mathrm dt=R^{-1}\int\mathcal E\,\mathrm dt=2NBS/R$。
因此
\[
\boxed{B=\frac{qR}{2NS}}
=\frac{4.5\times10^{-6}\times40}
{2\times60\times3.0\times10^{-6}}
=\boxed{0.50\,\mathrm T}.
\]
\excheck 电流计只记录翻转全过程累计电荷，不要求已知翻转速度。

\electroinductionproblem{300}{137}
图3.86边长为 $a$ 的正方形线框与长直导线电流 $I$ 共面，线框绕距长导线为 $b$ 的 OO' 轴转动 $180^\circ$；线框电阻 $R$、自感 $L$。求全过程流过线框的电荷。
\exsolution
直导线在距它 $r$ 处的场为
$B(r)=\mu_0I/(2\pi r)$。按图示轴线穿过正方形中心，前后有向通量分别为
$\Phi_i=(\mu_0Ia/2\pi)\ln[(b+a)/(b-a)]/2$ 与
$\Phi_f=-\Phi_i$（以图所用 $a$、$b$ 的长度约定），因而
$|\Delta\Phi|=(\mu_0Ia/2\pi)\ln[(b+a)/(b-a)]$。
回路满足
$L\dot I+RI=-\dot\Phi_{\rm ext}$。
初、末稳态回路电流均为零，积分得
$Rq=-\Delta\Phi$，故
\[
\boxed{|q|=\frac{\mu_0Ia}{2\pi R}
\ln\frac{b+a}{b-a}}.
\]
\excheck 本题全部过程累计电荷与自感 $L$ 无关；要求 $b>a$，否则线框会与直导线相交。

\electroinductionproblem{301}{138}
图3.87长直导线电流 $I_0$；位于其同一平面内的两平行导轨距直导线分别为 $a,b$，活动杆以恒速 $v$ 滑动，回路电阻 $R$。忽略其余电阻及自感，求 (a) 感应电流大小与方向；(b) 维持匀速所需外力。
\exsolution
磁场 $B(r)=\mu_0I_0/(2\pi r)$，活动杆从 $r=a$ 到 $r=b$，运动电动势
\[
|\mathcal E|=v\int_a^bB(r)\,\mathrm dr
=\frac{\mu_0I_0v}{2\pi}\ln\frac ba.
\]
因此
\[
\boxed{I_{\rm ind}=\frac{\mu_0I_0v}{2\pi R}\ln\frac ba}.
\]
由楞次定律感应电流在活动杆上产生阻碍其平动的磁力：
\[
F=I_{\rm ind}\int_a^bB(r)\,\mathrm dr
=\boxed{\frac vR\left(
\frac{\mu_0I_0}{2\pi}\ln\frac ba
\right)^2}.
\]
\excheck $Fv=I_{\rm ind}^2R$，严格满足机械功率转变为焦耳热。

\electroinductionproblem{302}{138}
图3.88一根质量 $m$、长 $\ell$ 的活动导体杆 AB 初速度为 $v_0$，在均匀磁场 $B$ 中沿无摩擦导轨滑动，导轨左端接电阻 $R$，忽略其他电阻及自感。求 (a) 直到停下的总位移；(b) 电阻产生的热量。
\exsolution
电流 $I=B\ell v/R$；安培阻力为
$F=-B\ell I=-B^2\ell^2v/R$，故
\[
m\dot v=-\frac{B^2\ell^2}{R}v,\qquad
v(t)=v_0e^{-t/\tau},\quad
\tau=\frac{mR}{B^2\ell^2}.
\]
总位移
\[
\boxed{s=\int_0^\infty v(t)\,\mathrm dt
=v_0\tau=\frac{mRv_0}{B^2\ell^2}}.
\]
动能最终全部变为焦耳热：
\[
\boxed{Q=\frac12mv_0^2}.
\]
\excheck 可直接验证
$\int_0^\infty I(t)^2R\,\mathrm dt
=\frac12mv_0^2$。

\electroinductionproblem{303}{138}
图3.89活动杆 AB 长 $\ell$、质量 $m$、电阻 $R$，在水平无摩擦导轨上受恒力 $F$ 向右运动，处于垂直导轨的均匀磁场 $B$ 中，初始静止。忽略自感和导轨电阻，求 $v(t)$。
\exsolution
运动中的磁感生电流为
$I=B\ell v/R$，阻尼安培力
$F_m=B^2\ell^2v/R$ 与运动反向。
运动方程
\[
m\dot v=F-\frac{B^2\ell^2}{R}v.
\]
设 $\gamma=B^2\ell^2/(mR)$，初值 $v(0)=0$，解得
\[
\boxed{v(t)=\frac{F}{m\gamma}
(1-e^{-\gamma t})
=\frac{FR}{B^2\ell^2}
\left(1-e^{-B^2\ell^2t/(mR)}\right)}.
\]
\excheck 初始加速度 $F/m$，终速 $FR/(B^2\ell^2)$。

\electroinductionproblem{304}{138}
图3.90(a)—(d)为位于纸面内的导线闭合结构，向纸内的均匀磁场开始减小。判断感生电流方向。
\exsolution
楞次定律：穿过任一有效闭合导电区域、方向\emph{指向纸内}的磁通正在减少，感生电流需产生\emph{同样指向纸内}的磁场，因此有效闭合环路电流沿顺时针。需进一步看连接结构：
(a) 圆周沿顺时针，穿过中心的直连导线因上下对称与两半环电势差相同，没有电流。
(b) 图3.90(b)的外闭合环顺时针。
(c) 两个同心闭合环均顺时针，径向短接导线中无电流。
(d) “8”字形两个环的有向边界相反，原图以左环顺时针为正，因此左环为顺时针，右环为逆时针。
\excheck 判断每一独立回路的实际拓扑，而不能把每一段导线都无条件标为顺时针；(a)(c)中的桥接支路因对称不传输电荷。
